%% LyX 2.4.4 created this file.  For more info, see https://www.lyx.org/.
%% Do not edit unless you really know what you are doing.
\documentclass[oneside,english]{book}
\usepackage[T1]{fontenc}
\usepackage[latin9]{inputenc}
\setcounter{secnumdepth}{3}
\setcounter{tocdepth}{3}
\usepackage{units}
\usepackage{amsmath}
\usepackage{amsthm}
\usepackage{amssymb}
\usepackage{cancel}
\usepackage{geometry}
\geometry{verbose,tmargin=1cm,bmargin=1cm,lmargin=1.5cm,rmargin=1.5cm}
\PassOptionsToPackage{normalem}{ulem}
\usepackage{ulem}

\makeatletter

%%%%%%%%%%%%%%%%%%%%%%%%%%%%%% LyX specific LaTeX commands.
%% Because html converters don't know tabularnewline
\providecommand{\tabularnewline}{\\}

%%%%%%%%%%%%%%%%%%%%%%%%%%%%%% Textclass specific LaTeX commands.
\theoremstyle{definition}
\newtheorem{defn}{\protect\definitionname}
\theoremstyle{definition}
\newtheorem{example}{\protect\examplename}
\ifx\proof\undefined\
  \newenvironment{proof}[1][\proofname]{\par
    \normalfont\topsep6\p@\@plus6\p@\relax
    \trivlist
    \itemindent\parindent
    \item[\hskip\labelsep
          \scshape
      #1]\ignorespaces
  }{%
    \endtrivlist\@endpefalse
  }
  \providecommand{\proofname}{Proof}
\fi
\theoremstyle{definition}
\newtheorem{xca}{\protect\exercisename}
\theoremstyle{plain}
\newtheorem{thm}{\protect\theoremname}
\theoremstyle{plain}
\newtheorem{cor}{\protect\corollaryname}
\theoremstyle{plain}
\newtheorem{lem}{\protect\lemmaname}

%%%%%%%%%%%%%%%%%%%%%%%%%%%%%% User specified LaTeX commands.
\usepackage{hyperref}
\usepackage[usenames,dvipsnames]{xcolor} %adds additional color options
\hypersetup{colorlinks=true, urlcolor=Blue}%Blue

\usepackage{multicol}

\usepackage{tikz}
\usepackage{tikz-3dplot}
\usetikzlibrary{calc,arrows,shapes,backgrounds,matrix,shapes.geometric,fit,trees,positioning}

\usetikzlibrary{patterns}
\usetikzlibrary{plotmarks}
\usepackage{pgfplots}
\pgfplotsset{compat=newest}

\usepackage{calc}
\usepackage[nomessages]{fp}

\usepackage{datetime}

%%%%%commands
\newcommand{\mb}[1]{$\mathbb{#1}$}
\newcommand{\funct}[2]{$f\colon#1\rightarrow#2$}
% Added by lyx2lyx
\setlength{\parskip}{\smallskipamount}
\setlength{\parindent}{0pt}

\makeatother

\usepackage{babel}
\providecommand{\corollaryname}{Corollary}
\providecommand{\definitionname}{Definition}
\providecommand{\examplename}{Example}
\providecommand{\exercisename}{Exercise}
\providecommand{\lemmaname}{Lemma}
\providecommand{\theoremname}{Theorem}

\begin{document}

\chapter{Groups}

We naturally think of many sets as having an associated operation.
Consider the familiar set of integers and the operation addition.
Adding integers gives another integer. Furthermore it does not matter
the order that integers are added. There is a special integer, $0$,
that does not affect other numbers when added to them. We can appreciate
that the set of integers under addition has a desirable structure.

We abstract these properties as \emph{group axioms}, and we will see
a highly diverse collection of structures with these same similar
properties. The ubiquity of groups in areas both within and outside
mathematics makes them a central organizing principle of mathematics.
The concept of group developed in the early nineteenth century in
the study of polynomial equations. 

First we formalize the above mentioned ``operation''.

\section{Groups}
\begin{defn}
Let $G$ be a set. A \emph{(binary) operation}\footnote{When parts of statements are given in parenthesis, it usually means
that future statements will leave that part out. So in the from here
on we will simply write ``operation'' as opposed to writing ``binary
operation''.} on the set $G$ is a function $f\colon G\times G\rightarrow G$.
\end{defn}
\begin{itemize}
\item $f$ is just a function sending two elements to a single element.
For example, define $f$ to be addition. $f(3,4)=3+4=7$.
\end{itemize}
\begin{example}
Define $f\colon\mathbb{R}\times\mathbb{R}\rightarrow\mathbb{R}$ as
$x+1$ for $x,y\in\mathbb{R}$.

\begin{itemize}
\item []$f(1,2)=3$ since $1+2=3$ and $f(2,-5)=2+(-5)=-3$ since, $2-5=-3$.
\end{itemize}
\end{example}

\begin{defn}
A \emph{group}, $G$, be a nonempty set with an associated (binary)
operation $\oplus$, that combines two elements to form another element
with the following properties for all $a,b,c\in G$.

\begin{enumerate}
\item Closure: $a\oplus b\in G$.
\item Associativity: $(a\oplus b)\oplus c=a\oplus(b\oplus c)$.
\item Identity: There exist an $0_{G}\in G$ such that $0_{G}\oplus a=a\oplus0_{G}=a$.
\item Inverse: There exists an $a^{-1}\in G$ such that $a\oplus(a^{-1})=0_{G}$.
\end{enumerate}
\end{defn}
\begin{itemize}
\item We use, $\oplus$, because the operation can be many more things than
$+$. Sometimes juxtaposition is used, e.g., $a\oplus b=ab$, in which
case it should not be confused with multiplication. 
\item The a set $G$ with associated operation $\oplus$ is often denoted
$(G,\oplus)$. For example the set of rational numbers under addition
might be written, $(\mathbb{Q},+)$.
\end{itemize}
\begin{example}
The set of rational numbers, $\mathbb{Q}$, under addition is a group.

\begin{proof}
$\mathbb{Q}$ certainly is nonempty; take $\frac{1}{2}\in\mathbb{Q}$
for instance. We can write any two rational numbers with a common
denominator. So WLOG we can write all rational numbers with the same
nonzero integer denominator, say $d$.

Closure: $\frac{x}{d}+\frac{y}{d}=\frac{x+y}{d}$. Since $x+y\in\mathbb{Z}$
it follows that $\frac{x+y}{d}\in\mathbb{Q}$.

Associativity: $\left(\frac{x}{d}+\frac{y}{d}\right)+\frac{z}{d}=\frac{x+y}{d}+\frac{z}{d}=\frac{x+y+z}{d}=\frac{x}{d}+\left(\frac{y}{d}+\frac{z}{d}\right)$.

Identity: $0=\frac{0}{d}$ and $\frac{0}{d}+\frac{x}{d}=\frac{x}{d}=\frac{x}{d}+\frac{0}{d}$

Inverse: $0$ is the denominator so for any $\frac{x}{d}$ we need
a $\frac{y}{d}$ such that $\frac{x}{d}+\frac{y}{d}=0$. Solve this
equation to find $y=-x$ so for any rational number $\frac{x}{d}$
, $\frac{-x}{d}$ is the additive inverse. 

Thus $\mathbb{Q}$ is a group.
\end{proof}
\end{example}
\begin{xca}
\label{ex:Z is a group under +}Show that the set of integers, $\mathbb{Z}$,
is a group under addition. 
\end{xca}

\begin{xca}
Prove that the set of even integers is a group under addition. 
\end{xca}

\begin{xca}
Prove that integers of $n$ multiples for $n\in\mathbb{Z}^{+}$ are
a group under addition.
\end{xca}

\begin{example}
$\left\{ 1,-1,i,-i\}\right\} $ where $i=\sqrt{-1}$ is a group under
multiplication. From the table,

\begin{tabular}{r|rrrr}
$\times$ & $1$ & $-1$ & $i$ & $-i$\tabularnewline
\hline 
$1$ & $1$ & $-1$ & $i$ & $-i$\tabularnewline
$-1$ & $-1$ & $1$ & $-i$ & $i$\tabularnewline
$i$ & $i$ & $-i$ & $-1$ & $1$\tabularnewline
$-i$ & $-i$ & $i$ & $1$ & $-1$\tabularnewline
\end{tabular}we can see that $1$ is the identity, every element has an inverse
($-i\cdot i=1$ for example), and the set is closed under multiplication.
The last property, associativity, follows by associativity of $\mathbb{C}$. 
\end{example}
To show that a set is not a group, it is usually easiest to find an
example where any one of the properties fails.
\begin{example}
$S=\left\{ 1,-1,i\right\} $ is not a group under multiplication. 
\end{example}
\begin{proof}
$-1(i)=-i$ and $-i\notin S$ so it is not closed. Alternatively,
$i\cdot x=1\Rightarrow x=-i$. So the inverse of $i$ is $-i$, $-i\notin S$. 
\end{proof}
\begin{xca}
\label{matrices are a group}Prove that the set of all $2\times2$
matrices with determinant $1$ and entries from $\mathbb{R}$ is a
group under matrix multiplication. 
\end{xca}

\begin{xca}
Show that the following are \emph{not} groups under the given operation.

\begin{enumerate}
\item $S=\left\{ z\colon z\in\mathbb{Z}\mbox{ and }z=\sqrt{5}\right\} $under
addition: $(S,+)$
\item $\mathbb{Q}$, the set of rationals under multiplication: $(\mathbb{Q},\times)$
\item The set of odd integers under addition.
\item $\mathbb{Z}$ under multiplication: $(\mathbb{Z},\times)$
\end{enumerate}
\end{xca}

\begin{xca}
Let $S=\left\{ a+b\sqrt{5}|\,a,b\mbox{ are not both zero and }a,b\in\mathbb{Q}\right\} $.
Prove that $S$ is a group. (Hint: rationalize the denominator to
show that the inverse of $a+b\sqrt{5}$ is in $S$.)
\end{xca}

\begin{xca}
Show that $\left\{ 1,2,3\right\} $ under multiplication modulo $4$
is \emph{not }a group but that $\left\{ 1,2,3,4\right\} $ under multiplication
modulo $5$ is a group. (Example: $2\cdot3=6$ and $6=2\bmod4$ so
$[6]=[2]$ and thus $2\cdot3$ \emph{is} in $\left\{ 1,2,3\right\} $). 
\end{xca}
It is easy to take the next few theorems for granted. Especially the
cancellation law. You probably have been using it for some time without
giving it much thought. However we cannot assume such properties until
it is proven. Imagine trying to solve a simple linear equation without
it.
\begin{thm}
\textbf{\emph{\label{thm:Cancellation-Law.}Cancellation Law. }}Right
and left cancellation hold in a group, i.e., $b\oplus a=c\oplus a\iff b=c$
and $a\oplus b=a\oplus c\iff b=c$ for any elements $a,b,c$ in the
group.

\begin{proof}
Let $G$ be group with elements $a,b,c$. The inverses of the elements,
$a^{-1},b^{-1},c^{-1}$, are also in $G$ as $G$ is a group. First
we show the case for the right. Suppose that $b\oplus a=c\oplus a$
then 
\begin{eqnarray*}
(b\oplus a)\oplus a^{-1} & = & (c\oplus a)\oplus a^{-1}\\
\Rightarrow b\oplus(a\oplus a^{-1}) & = & c\oplus(a\oplus a^{-1})\mbox{ by associativity}\\
\Rightarrow b\oplus0_{G} & = & c\oplus0_{G}\mbox{ since }a\oplus a^{-1}=0_{G}\\
b & = & c
\end{eqnarray*}

Similarly we can show the case for the left. That $a\oplus b=a\oplus c\iff b=c$. 
\end{proof}
\end{thm}

\begin{thm}
\label{thm:identity-is-unique}A group has one and only one identity
element, i.e., a group's identity element is unique.

\begin{proof}
Let $G$ be a group. Since $G$ is a group it has an identity element
$0_{G}$. Now suppose there is some other $e\in G$ such that $e\oplus a=a\oplus e=a$
for any $a\in G$. We need to show that $e=0_{G}$. Since $G$ is
a group $a$ has an inverse $a^{-1}$. 
\begin{eqnarray*}
(0_{G}\oplus a)\oplus a^{-1} & = & (e\oplus a)\oplus a^{-1}\mbox{ since }o_{G}\oplus a=e\oplus a=a\\
0_{G}\oplus(a\oplus a^{-1}) & = & e\oplus(a\oplus a^{-1})\mbox{ by associativity}\\
0_{G}\oplus0_{G} & = & e\oplus0_{G}\mbox{ since }a\oplus a^{-1}=0_{G}\\
0_{G} & = & e
\end{eqnarray*}
\end{proof}
\end{thm}

\begin{thm}
\label{thm:Inverses-are-unique}An element in a group has one and
only one inverse. 

\begin{proof}
Let $0_{G}$ be the unique identity of a group $G$. Suppose there
exists an $a^{-1},x\in G$ such that $a\oplus a^{-1}=a\oplus x=0_{G}$.
Then 
\begin{eqnarray*}
x & = & x\oplus0_{G}\\
 & = & x\oplus(a\oplus a^{-1})\\
 & = & (x\oplus a)\oplus a^{-1}\\
 & = & 0_{G}\oplus a^{-1}=a^{-1}
\end{eqnarray*}
So $x=a^{-1}$ as needed. 
\end{proof}
\end{thm}
\begin{itemize}
\item Note that since the identity of a group is its own inverse, i.e.,
$0_{G}\oplus0_{G}=0_{G}$, theorem \ref{thm:identity-is-unique} can
be proved as a corollary of theorem \ref{thm:Inverses-are-unique}. 
\end{itemize}
\begin{xca}
$\bigstar$ Addition in $\mathbb{Z}_{n}$ is defined as $\left[a\right]+\left[b\right]=\left[a+b\right]$.
Multiplication in $\mathbb{Z}_{n}$ is defined as $\left[a\right]\cdot\left[b\right]=\left[a\cdot b\right]$.
For example, $\mathbb{Z}_{4}=\left\{ \left[0\right],\left[1\right],\left[2\right],\left[3\right]\right\} $
for convenience we write $\mathbb{Z}_{4}=\left\{ 0,1,2,3\right\} $.
$2+3=5$ and $5=1\bmod4$ so $\left[5\right]=\left[1\right]$. Show
that $\mathbb{Z}_{2}\times\mathbb{Z}_{3}$ is a group under modular
addition. 
\end{xca}
\begin{defn}
A group is \emph{abelian}\footnote{Named in honor of the mathematician Niels Henrik Abel\emph{. }Abel
made many outstanding contributions to a variety of fields; though
he died at the age of 26. Abelian groups are also called commutative
groups.}\emph{ }if $a\oplus b=b\oplus a$ for all $a,b$ in the group.
\end{defn}
\begin{example}
$\mathbb{Z}$ is an abelian group under addition.
\end{example}

\begin{xca}
Prove that the set of all $n\times n$ matrices is an abelian group
under \uline{addition}.
\end{xca}

\begin{xca}
Prove that the set of all $2\times2$ matrices with determinant $1$
and entries from $\mathbb{R}$ is a non-abelian group under matrix
multiplication (we already know that it is a group by exercise \ref{matrices are a group}. 
\end{xca}

\begin{xca}
$\bigstar$ Consider a square placed on the Cartesian plane such that
the square's center is at the origin and its corners are on the axes
(like a diamond; see the figure below). Any motion of the square which
keeps the center at the origin and corners on the axes has the same
result as one of eight motions: $r_{0}\equiv\mbox{ rotation of }0^{\circ}$,
$r_{1}\equiv\mbox{ rotation of }90^{\circ}$, $r_{3}\equiv\mbox{ rotation of }180^{\circ}$,
$d_{x}\equiv\mbox{ }x-\mbox{axis reflection}$, $h_{+}\equiv\,y=x\mbox{ reflection}$,
or $h_{-}\equiv y=-x\mbox{ reflection}$. \\
\hspace*{.5cm}\newcommand{\xmax}{1.5}
\newcommand{\xmin}{-1.5}
\newcommand{\xstep}{0} %must be xmin+n
\newcommand{\ymax}{1.5}
\newcommand{\ymin}{-1.5}
\newcommand{\ystep}{0} %must be ymin+n
%%%%%%%
\begin{tikzpicture} 
\begin{axis}[height=3.5cm, 
	width=3.5cm, 
	axis lines=middle,
	x axis line style={-},y axis line style={-}, 
	grid=none,
	%axis line style={-latex}, 
	xmin=\xmin,xmax=\xmax, 
	ymin=\ymin,ymax=\ymax,
	%restrict y to domain=-1:10,
	no markers, 
	xlabel={$$},ylabel={$$},
	ticks=none,
	]
	
	\addplot[very thick,blue] coordinates{(1,0) (0,1)};
	\addplot[very thick,blue] coordinates{(0,1) (-1,0)};
	\addplot[very thick,blue] coordinates{(-1,0) (0,-1)};
	\addplot[very thick,blue] coordinates{(0,-1) (1,0)};

	\node at (axis cs:1.3,.3) {$1$};
	\node at (axis cs:-.3,1.3) {$2$};
	\node at (axis cs:-1.3,-.3) {$3$};
	\node at (axis cs:.3,-1.3) {$4$};
	
\end{axis}
\end{tikzpicture}\\
Show that $D_{4}=\left\{ r_{0},r_{1},r_{3},d_{x},d_{y},h_{+},h_{-}\right\} $
is a non-abelian group under composition of these motions. The set
$D_{4}$ is called the \emph{dihedral group of degree 4}. 
\end{xca}

\begin{thm}
Let $G$ be a group and $a,b\in G$ then $(a\oplus b)^{-1}=b^{-1}\oplus a^{-1}$.

\begin{proof}
The inverse of $(a\oplus b)^{-1}$ is $(a\oplus b)$. Now $(b^{-1}\oplus a^{-1})\oplus(a\oplus b)=b^{-1}\oplus(a^{-1}\oplus a)\oplus b=b^{-1}\oplus0_{G}\oplus b=b^{-1}\oplus b=0_{G}$.
So then $b^{-1}\oplus a^{-1}$ is the inverse of $(a\oplus b)$, and
by theorem \ref{thm:identity-is-unique} 
\end{proof}
\end{thm}
\begin{xca}
Prove that $G$ is an Abelian group if and only is $\left(a\oplus b\right)^{-1}=a^{-1}\oplus b^{-1}$
for all $a,b\in G$. 
\end{xca}

\begin{xca}
Suppose that if $G$ is a group such that if $a,b,c\in G$ and $a\oplus b=c\oplus a$
then $b=c$. Prove that $G$ is an Abelian group, i.e., that $a\oplus b=b\oplus a$
for all $a,b\in G$. Hint: consider the element $a\oplus(b\oplus a)=(a\oplus b)\oplus a$.
\end{xca}

\begin{xca}
$\bigstar$ Prove that if in a group $G$ that $\left(a\oplus b\right)^{2}=a^{2}\oplus b^{2}$
then $G$ is Abelian.
\end{xca}


\section{Subgroups}
\begin{defn}
If a\emph{ }subset $H$ of a group $G$ is a group under the operation
of $G$, then $H$ is a \emph{subgroup }of $G$, denoted $H\leq G$.
\end{defn}
\begin{itemize}
\item Because $H$ is a group, it is implied in this definition that all
subgroups are non-empty.
\end{itemize}
\begin{example}
\label{ex: 2+Z10 is a subgrop of Z10}$H=\left\{ 2,4,6,8,0\right\} $
is a subgroup of the group $\mathbb{Z}_{10}$ under modular addition.

\begin{proof}
From example \ref{ex:Z is a group under +}, we know that $\mathbb{Z}_{10}$
is a group under addition modulo $10$. \textbf{$H\neq\emptyset$}
and $H\subseteq\mathbb{Z}_{10}$. Using the Cayley table,

\begin{center}{\bfseries{}%
\begin{tabular}{c|ccccc}
$\oplus$ & 0 & 2 & 4 & 6 & 8\tabularnewline
\hline 
0 & 0 & 2 & 4 & 6 & 8\tabularnewline
2 & 2 & 4 & 6 & 8 & 0\tabularnewline
4 & 4 & 6 & 8 & 0 & 2\tabularnewline
6 & 6 & 8 & 0 & 2 & 4\tabularnewline
8 & 8 & 0 & 2 & 4 & 6\tabularnewline
\end{tabular}}\end{center}ee can see that it is closed, has identity element $0$
(the same as $\mathbb{Z}_{10}$), and each element has an inverse
($6+4\equiv0\bmod10$ for example). Associativity holds in modular
addition; so $H$ is a group , and hence a subgroup of $\mathbb{Z}_{10}$.
Written $H\leq\mathbb{Z}_{10}$. 
\end{proof}
\end{example}
A subset $H$ of a group $G$ will always have the associativity property
of $G$. So we need only check for closure, existence of an identity
and inverses. The relation between the identity and inverses allows
us check these properties in a single steo with the following theorem.
\begin{thm}
\textbf{\emph{\label{thm:Subgroup-Test.}Subgroup Test. }}Let $G$
be a group and $H$ be a nonempty subset of $G$. Then $H$ is a subgroup
of $G$ if and only if $a\oplus b^{-1}\in H$ for all $a,b\in H$.

\begin{proof}
$(\Rightarrow)$ Suppose that $H\leq G$, and let $a,b\in H$. $H$
is a group every so element has an inverse, i.e., there exists a $b^{-1}\in H$.
Groups are also closed so $a\oplus b^{-1}\in H$.

$(\Leftarrow)$ Suppose that $a\oplus b^{-1}\in H$, and let $a,b\in H$
for some nonempty subset of $G$. 

\begin{enumerate}
\item []Associativity: The operation, $\oplus$, in $H$ is the same as
that of $G$ so $\oplus$ is associative in $H$. 
\item []Identity: By the hypothesis for any $a\in H\Rightarrow a\oplus a^{-1}=0_{G}\in H$.
So $H$ has the same identity element as $G$ (at least one such element
exists since $H\neq\emptyset$).
\item []Inverse: $0_{G}\in H$ (shown above) so for any $a\in H$ we have
that $0\oplus a^{-1}=a^{-1}\in H$.
\item []Closure: We just showed that for and $a,b\in H\Rightarrow a^{-1},b^{-1}\in H$.
$a\oplus b=a\oplus(b^{-1})^{-1}\in H$.
\end{enumerate}
So $H\leq G$.
\end{proof}
\end{thm}
\begin{itemize}
\item Very often, the operation is just addition. Writing $+$ instead of
$\oplus$, the subgroup test is then: 

\begin{itemize}
\item []A nonempty subset $H$ of a group $G$ is a subgroup of $G$ if
and only if $a-b\in H$ for all $a,b\in H$.
\end{itemize}
\end{itemize}
\begin{example}
$H=\left\{ 2,4,6,8,0\right\} $ is a subgroup of the group $\mathbb{Z}_{10}$
under modular addition. 

\begin{proof}
For any $a,b\in H$, $b^{-1}=-b$ since $b-b=0$. Using the subgroup
test, $a-b\in G$ since $\mathbb{Z}_{10}$ is a group. So we need
only show that $2|a-b$. $2|a$ and $2|b\Rightarrow a=2n$ and $b=2m$
for some $m,n\in\mathbb{Z}$, and $2|2(n-m)\Rightarrow2|a-b$. So
then $a-b\in H$. 
\end{proof}
\end{example}

\begin{example}
Show that $\mathbb{Q}$ is a subgroup of the group $\mathbb{R}$ under
addition. 

\begin{proof}
$\mathbb{Q}$ is a nonempty subset of $\mathbb{R}$. Let $\frac{a}{b},\frac{c}{d}\in\mathbb{Q}$.
$\frac{c}{d}^{-1}=-\frac{c}{d}$ since $\frac{c}{d}-\frac{c}{d}=0$.
$\frac{a}{b}-\frac{c}{d}=\frac{ad-bc}{bd}\in\mathbb{Q}$. So then
$\mathbb{Q}\leq\mathbb{R}$ by the subgroup test. 
\end{proof}
\end{example}
\begin{xca}
Let $\left(\mathbb{Q},+\right)$ be the group of rational numbers
under addition. 

\begin{enumerate}
\item Is $\mathbb{Z}$ a subgroup of $\left(\mathbb{Q},+\right)$?
\item Is $\mathbb{Z}^{+}=\left\{ 0,1,2,3,\dots\right\} $ a subgroup of
$\left(\mathbb{Q},+\right)$?
\end{enumerate}
\end{xca}

\begin{xca}
Determine if the following are subgroups.
\end{xca}
\begin{enumerate}
\item Is $\mathbb{Q}$ a subgroup of the group, $\left(\mathbb{C},+\right)$,
the complex numbers under addition?
\item Is $\mathbb{Z}-\left\{ 0\right\} $ a subgroup of the group, $\left(\mathbb{Q}-\left\{ 0\right\} ,\cdot\right)$,
the non-zero rational numbers under multiplication?
\end{enumerate}

\begin{xca}
Prove that the intersection of two subgroups $H$ and $K$ of a group
$G$ is a subgroup.
\end{xca}

\begin{xca}
Let $H_{1},H_{2},\dots$ be subgroups of $G$. Prove that $H_{1}\cap H_{2}\cap\cdots$
is a subgroup of $G$.
\end{xca}

\begin{xca}
Recall that $\mathbb{Z}_{6}=\left\{ 0,1,2,3,4,5\right\} $ is a group
under modular addition. 

\begin{enumerate}
\item Show that $H=\left\{ 0,2,4\right\} $ and $K=\left\{ 0,3\right\} $
are subgroups of $\mathbb{Z}_{6}$.
\item Show that $H\cup K$ is \emph{not }a subgroup of $\mathbb{Z}_{6}$.
\item Is the union of two subgroups of $G$ necessarily a subgroup $G$?
\end{enumerate}
\end{xca}

\begin{xca}
$\bigstar$ Let $\mathbb{Z}_{6}=\left\{ 0,1,2,3,4,5\right\} $ and
$\mathbb{Z}_{4}=\left\{ 0,1,2,3\right\} $. $\mathbb{Z}_{6}\times\mathbb{Z}_{4}$
is a group under modular addition. Show that $H=\left\{ \left(0,0\right),\left(3,0\right),\left(0,2\right),\left(3,2\right)\right\} $
is a subgroup. 
\end{xca}

\begin{xca}
Let $G$ be an abelian group and $H$ be a subgroup of $G$. Prove
that $S(H)=\left\{ x|x\in G\mbox{ and }x\cdot x\in H\right\} $ is
also a subgroup of $G$.
\end{xca}

\begin{xca}
Let $G$ be a group. Prove that $G$ has a subgroup with one element.
This group is called the \emph{trivial group}.
\end{xca}

\section{Finite Groups and Cyclic Groups}

Groups with finitely many elements have interesting arithmetic properties. 
\begin{defn}
The number of elements in a group its \emph{order}, denoted $\left|G\right|$. 
\end{defn}
This is exactly how we defined the order of a set. For example, $\left|\mathbb{Z}_{10}\right|=10$
and from example \ref{ex: 2+Z10 is a subgrop of Z10}, $\left|H\right|=5$. 

\begin{defn}
The \emph{order of an element }$a$ in a group $G$ is the smallest
positive integer $n$ such that $a^{n}=0_{G}$ ($a^{n}=\underset{n\mbox{ times}}{\underbrace{a\oplus a\oplus\dots\oplus a}}).$
\end{defn}
\begin{example}
From example \ref{ex: 2+Z10 is a subgrop of Z10}, $2+2+2+2+2=10\equiv0\bmod10$.
So in $H$, $\left|2\right|=5$. In $\mathbb{Z}_{10}$, $\left|2\right|=5$
and $\left|1\right|=10$. 
\end{example}
\begin{defn}
The set \emph{generated by $a$ }is the set $\left\langle a\right\rangle =\left\{ a^{n}\colon n\in\mathbb{Z}\right\} $.
The element $a$ is called the \emph{generator }of $\left\langle a\right\rangle $.
\end{defn}
\begin{itemize}
\item If $a$ is an element of a group $G$ with $\left|a\right|=n$ then
$a^{k}\in G$ for any integer $k$, and $\left\langle a\right\rangle =\left\{ 0_{G},a,a^{2},\dots,a^{n-1}\right\} $.
\end{itemize}
\begin{example}
Let $2\in\mathbb{Z}_{10}$. $\left\langle 2\right\rangle =\left\{ 2,3,6,8,0\right\} $.
Recall that we have already showed that this is a subgroup of $\mathbb{Z}_{10}$.
Since $\left\langle 2\right\rangle =\left\{ 2,3,6,8,0\right\} $ we
say that $2$ generates the group $\left\{ 2,3,6,8,0\right\} $.
\end{example}

\begin{example}
Find all generators of $\mathbb{Z}_{6}$.

\begin{itemize}
\item []$\left\langle 1\right\rangle =\left\{ 1,2,3,4,5,0\right\} $(note:
$1$ is a generator of $\mathbb{Z}_{n}$ for any $n$) and $\left\langle 5\right\rangle =\left\{ 5,4,3,2,1,0\right\} $
so $\left\langle 1\right\rangle =\left\langle 5\right\rangle =\mathbb{Z}_{6}$.
\end{itemize}
\end{example}
\begin{xca}
Find all generators of $\mathbb{Z}_{8}$ and $\mathbb{Z}_{10}$.
\end{xca}

Some of your favorite sets are generated by a single element. For
example, any integer $n$ can be written as $\pm1\pm1\pm\cdots\pm1=n$.
So $1$ generates $\mathbb{Z}$. Similarly $2$ and $3$ generate
$2\mathbb{Z}$ and $3\mathbb{Z}$ respectively. 
\begin{defn}
Let $G$ be a group. If there is a single element $g\in G$ such that
every element in $G$ is equal to $g^{n}=\underset{n\mbox{ times}}{\underbrace{g\oplus g\oplus\dots\oplus g}}$
for some $n\in\mbox{\ensuremath{\mathbb{Z}}}$ then $G$ is a \emph{cyclic
}group and \emph{$g$ }is called the \emph{generator} of $G$, and
we write $G=\left\langle g\right\rangle $.
\end{defn}
\begin{itemize}
\item $\left\langle g\right\rangle =\left\{ g^{n}\colon n\in\mathbb{Z}\right\} $
\end{itemize}
\begin{xca}
Prove that $\mathbb{Z}$ is a cyclic group. 
\end{xca}

\begin{xca}
Prove that $\mathbb{Z}_{n}$ is a cyclic group under modular addition
for any $n\in\mathbb{Z}$ such that $n>0$. (you can use the fact
that it is a group).
\end{xca}
\begin{thm}
\label{thm:All-cyclic-groups-are-abelian}All cyclic groups are abelian.

\begin{proof}
Let $G=\left\langle a\right\rangle $ be a group. So any $x,y\in G$
can be written as $a^{i}$ and $a^{j}$ respectively. $a^{i}\oplus a^{j}=a^{i+j}=a^{j+i}=a^{i}\oplus a^{j}$.
\end{proof}
\end{thm}

\begin{thm}
\label{thm:<a> is a subgroup of G}If $G$ is a group and $a\in G$
then $\left\langle a\right\rangle \leq G$.

\begin{proof}
$a\in\left\langle a\right\rangle $ so $\left\langle a\right\rangle $
is nonempty, and $G$ is closed so $a^{m}\in G$ for any integer $m$.
So $\left\langle a\right\rangle $ is a nonempty subset of $G$. Now
we use the subgroup test (theorem \ref{thm:Subgroup-Test.}). 

Let $a^{i}$ and $a^{j}$ be two elements in $\left\langle a\right\rangle $.
$a^{i}\oplus a^{-j}=a^{i-j}\in\left\langle a\right\rangle $ since
$i-j\in\mathbb{Z}$. 
\end{proof}
\end{thm}
\begin{itemize}
\item Theorem \ref{thm:All-cyclic-groups-are-abelian} gives us another
way to show that a group is abelian. 
\item Theorem \ref{thm:<a> is a subgroup of G} gives us an easy way to
find abelian subgroups of any group. 
\end{itemize}
\begin{example}
$\left\langle 2\right\rangle ,\left\langle 3\right\rangle ,\dots,\left\langle n\right\rangle $
are all abelian subgroups of $\mathbb{Z}$. 
\end{example}

\begin{thm}
\label{thm:a^i=00003Da^j iff n|i-j}Let $G$ be a group and $a\in G$.
If $\left|a\right|=n$ then $a^{i}=a^{j}$ if and only if $n|i-j$.

\begin{proof}
$(\Leftarrow)$ Let $n|i-j$. Then $i-j=nq$ and 
\[
a^{i-j}=a^{nq}=\left(a^{n}\right)^{q}=(0_{G})^{q}=0_{g}
\]
as needed.

$(\Rightarrow)$ $a^{i}=a^{j}$ implies that $a^{i}\oplus a^{-j}=a^{i-j}=0_{G}$.
By the division algorithm we can write $i-j=qn+r$ for $q,r\in\mathbb{Z}$
and with $0\leq r<n$. 
\[
0_{G}=a^{i-j}=a^{qn+r}=(a^{n})^{q}\oplus a^{r}=(0_{G})^{q}\oplus a^{r}=a^{r}=0_{G}
\]
Thus $r=0$, and $i-j=qn\Rightarrow n|i-j$ as needed.
\end{proof}
\end{thm}
\begin{itemize}
\item An immediate consequence of this theorem is that the order of the
group generated by $a$ is the same as the order of $a$, i.e., $\left|a\right|=\left|\left\langle a\right\rangle \right|$.
\end{itemize}
\begin{example}
Consider $\mathbb{Z}_{30}$ a group under addition modulo $30$. $\mathbb{Z}_{30}$
has identity $0$ and is generated by $1$. $\left|5\right|=6$ since
\[
5+5+5+5+5+5=(6)5\equiv0\bmod30
\]

What is $5^{40}$ $\bmod30$? $5^{4}=4(5)=20$ and $6|40-4$ so then
$5^{40}\equiv20\bmod30$ as well. More generally, $6|x-4\iff x=6q+4$
for some $q\in\mathbb{Z}$. So then $5^{6q+4}\equiv20\bmod30$. 
\end{example}

\begin{cor}
\label{cor: a^k=00003D0 iff n|k}Let $G$ be a group and $a\in G$.
$\left|a\right|=n$ and $a^{k}=0_{G}$ if and only if $n|k$.

\begin{proof}
$a^{k}=a^{k-0}=0_{G}$ we know by the theorem above that $n|k-0$
so $n|k$. The other direction also follows from the theorem.
\end{proof}
\end{cor}
\begin{figure}
\center\begin{tikzpicture}
	\def \n {5} 
	\def \radius {2.5cm} 
	\def \margin {8} % margin in angles, depends on the radius
	\foreach \s in {1,...,\n} {   
		\node[draw,circle] at ({360/\n * (\s - 1)}:
		\radius) {$\s$};   \draw[->, >=latex] ({360/\n * (\s - 1)+\margin}:\radius) arc ({360/\n * (\s - 1)+\margin}:{360/\n * (\s)-\margin}:\radius);
		} 
\end{tikzpicture}

\caption{A graph representing the cyclic group $\mathbb{Z}_{6}$.}
\end{figure}


\section{Permutation Groups}

Permutation groups are functions from a set $A$ back onto $A$.\footnote{Permutations were the first groups studied by mathematicians. In 1832
the French mathematician �variste Galois gave a criterion for the
solvability of some polynomial equations in terms of the symmetry
group of its roots (solutions), and the elements of these groups correspond
to certain permutations of the roots. Arthur Cayley's work on permutations
in 1854 gave the first abstract definition of a finite group. } They can be thought of ways of ``rearranging'' the elements of
a set.
\begin{defn}
A \emph{permutation }of a set $A$ is a function from $A$ to $A$
that is one-to-one and onto (a bijection). A \emph{permutation group
}of a set $A$ is a set of permutations of $A$ that is a group under
function composition.
\end{defn}
\begin{example}
\label{ex: simple permutations }Let $\sigma\colon\left\{ 1,2,3\right\} \rightarrow\left\{ 1,2,3\right\} $
and $\gamma\colon\left\{ 1,2,3\right\} \rightarrow\left\{ 1,2,3\right\} $
be two bijections from $\left\{ 1,2,3\right\} $ to $\left\{ 1,2,3\right\} $
such that 

\begin{tabular}{rlc}
$\sigma(1)=2$ &  & $\gamma(1)=2$\tabularnewline
$\sigma(2)=3$ & and & $\gamma(2)=1$\tabularnewline
$\sigma(3)=1$ &  & $\gamma(3)=3$\tabularnewline
\end{tabular}

Function composition is then well defined, and $\sigma\circ\gamma\colon\{1,2,3\}\rightarrow\{1,2,3\}$
is also a bijective function from $\{1,2,3\}$ to $\{1,2,3\}$. For
example $\gamma(\sigma(1))=\gamma(2)=1$.
\end{example}
There are different ways to represent permutations on a set $A$.
As permutations map every element in $A$ to one and only one element,
the mapping of any element forms a cycle $a\rightarrow b\rightarrow c\rightarrow\dots\rightarrow a$.
This cycle may not include all the elements in $A$. Indeed it may
just be $a\rightarrow a$, but it will always end where it started.
So a map that takes $a\rightarrow b\rightarrow c\rightarrow\dots\rightarrow a$
can be written conveniently as $(a,b,c,\dots,a)$ and $a\rightarrow a$
as simply $a$. Writing permutations in this way is called \emph{cyclic
form}. 
\begin{example}
Consider the permutations (from the example \ref{ex: simple permutations }
above) $\sigma\colon\left\{ 1,2,3\right\} \rightarrow\left\{ 1,2,3\right\} $
and $\gamma\colon\left\{ 1,2,3\right\} \rightarrow\left\{ 1,2,3\right\} $
be two bijections from $\left\{ 1,2,3\right\} $ to $\left\{ 1,2,3\right\} $
such that 

\begin{tabular}{rlc}
$\sigma(1)=2$ &  & $\gamma(1)=2$\tabularnewline
$\sigma(2)=3$ & and & $\gamma(2)=1$\tabularnewline
$\sigma(3)=1$ &  & $\gamma(3)=3$\tabularnewline
\end{tabular}

Graphs representing $\sigma$ and $\gamma$ are shown in figures \ref{fig:A-graph-representing (1,2,3)}
and \ref{fig:Graphs-representing-(1,2)(3)} respectively.
\end{example}
\begin{figure}[h]
\center\begin{tikzpicture}[scale=.75]
	\def \n {3} %number of elements
	\def \x {3} %endnumber
	\def \ftn {$\sigma$} %function
	\def \radius {2cm}
	\def \radiusNode {2.25cm} %places nodes
	\def \label {$(1,2,3)$} %center label
	\def \margin {8} % margin in angles, depends on the radius
	\foreach [count=\i] \s in {1,2,\x} {   
		\node at ({360/\n * (\i - 1)}:
			\radius) {$\s$};
		\node at ({360/\n * (\i - 1) + 360/(2*\n)}:
			\radiusNode) {\ftn};
		\draw[->, >=latex] ({360/\n * (\i - 1)+\margin}:\radius) arc ({360/\n * (\i - 1)+\margin}:{360/\n * (\i)-\margin}:\radius);
		}
	\node {\label};
\end{tikzpicture}

\caption{\label{fig:A-graph-representing (1,2,3)}A graph representing the
permutation of $\sigma=(1,2,3)$ on $\{1,2,3\}$.}
\end{figure}

\begin{figure}[h]
\begin{minipage}[t]{0.5\textwidth}%
\begin{center}
\begin{tikzpicture}[scale=.75]
	\def \n {2} %number of elements
	\def \x {2} %endnumber
	\def \ftn {$\sigma$} %function
	\def \radius {2cm}
	\def \radiusNode {2.25cm} %places nodes
	\def \label {$(1,2)$} %center label
	\def \margin {8} % margin in angles, depends on the radius
	\foreach [count=\i] \s in {1,\x} {   
		\node at ({360/\n * (\i - 1)}:
			\radius) {$\s$};
		\node at ({360/\n * (\i - 1) + 360/(2*\n)}:
			\radiusNode) {\ftn};
		\draw[<-, >=latex] ({360/\n * (\i - 1)+\margin}:\radius) arc ({360/\n * (\i - 1)+\margin}:{360/\n * (\i)-\margin}:\radius);
	}
	\node {\label};
\end{tikzpicture}
\par\end{center}%
\end{minipage}%
\begin{minipage}[t]{0.5\textwidth}%
\begin{center}
\begin{tikzpicture}[scale=.75]
	\def \n {1} %number of elements
	\def \x {3} %endnumber
	\def \ftn {$\sigma$} %function
	\def \radius {2cm}
	\def \radiusNode {2.25cm} %places nodes
	\def \label {$(3)$} %center label
	\def \margin {8} % margin in angles, depends on the radius
	\foreach [count=\i] \s in {\x} {   
		\node at ({360/\n * (\i - 1)}:
			\radius) {$\s$};
		\node at ({360/\n * (\i - 1) + 360/(2*\n)}:
			\radiusNode) {\ftn};
		\draw[<-, >=latex] ({360/\n * (\i - 1)+\margin}:\radius) arc ({360/\n * (\i - 1)+\margin}:{360/\n * (\i)-\margin}:\radius);
		}
	\node {\label};	
\end{tikzpicture}
\par\end{center}%
\end{minipage}

\caption{\label{fig:Graphs-representing-(1,2)(3)}Graphs representing the permutation
of $\gamma=(1,2)(3)$ on $\{1,2,3\}$.}
\end{figure}

\begin{example}
Let $\alpha$ be a permutation on $\{1,2,3,4,5,6\}$ defined as follows: 

\begin{tabular}{rlc}
$\alpha(1)=2$ &  & $\alpha(4)=3$\tabularnewline
$\alpha(2)=1$ & and & $\alpha(5)=5$\tabularnewline
$\alpha(3)=6$ &  & $\alpha(6)=4$\tabularnewline
\end{tabular}

In cycle form $\alpha=(1,2)(3,6,4)(5)$. See the graph in figure 
\begin{figure}[h]
\begin{minipage}[t]{0.33\textwidth}%
\begin{center}
\begin{tikzpicture}[scale=.75]
	\def \n {2} %number of elements
	\def \x {2} %endnumber
	\def \ftn {$\sigma$} %function
	\def \radius {2cm}
	\def \radiusNode {2.25cm} %places nodes
	\def \label {$(1,2)$} %center label
	\def \margin {8} % margin in angles, depends on the radius
	\foreach [count=\i] \s in {1,\x} {   
		\node at ({360/\n * (\i - 1)}:
			\radius) {$\s$};
		\node at ({360/\n * (\i - 1) + 360/(2*\n)}:
			\radiusNode) {\ftn};
		\draw[<-, >=latex] ({360/\n * (\i - 1)+\margin}:\radius) arc ({360/\n * (\i - 1)+\margin}:{360/\n * (\i)-\margin}:\radius);
	}
	\node {\label};
\end{tikzpicture}
\par\end{center}%
\end{minipage}%
\begin{minipage}[t]{0.33\textwidth}%
\begin{center}
\begin{tikzpicture}[scale=.75]
	\def \n {3} %number of elements
	\def \x {4} %endnumber
	\def \ftn {$\sigma$} %function
	\def \radius {2cm}
	\def \radiusNode {2.25cm} %places nodes
	\def \label {$(3,6,4)$} %center label
	\def \margin {8} % margin in angles, depends on the radius
	\foreach [count=\i] \s in {3,6,\x} {   
		\node at ({360/\n * (\i - 1)}:
			\radius) {$\s$};
		\node at ({360/\n * (\i - 1) + 360/(2*\n)}:
			\radiusNode) {\ftn};
		\draw[->, >=latex] ({360/\n * (\i - 1)+\margin}:\radius) arc ({360/\n * (\i - 1)+\margin}:{360/\n * (\i)-\margin}:\radius);
		}
	\node {\label};	
\end{tikzpicture}
\par\end{center}%
\end{minipage}%
\begin{minipage}[t]{0.33\textwidth}%
\begin{center}
\begin{tikzpicture}[scale=.75]
	\def \n {1} %number of elements
	\def \x {5} %endnumber
	\def \ftn {$\sigma$} %function
	\def \radius {2cm}
	\def \radiusNode {2.25cm} %places nodes
	\def \label {$(5)$} %center label
	\def \margin {8} % margin in angles, depends on the radius
	\foreach [count=\i] \s in {\x} {   
		\node at ({360/\n * (\i - 1)}:
			\radius) {$\s$};
		\node at ({360/\n * (\i - 1) + 360/(2*\n)}:
			\radiusNode) {\ftn};
		\draw[<-, >=latex] ({360/\n * (\i - 1)+\margin}:\radius) arc ({360/\n * (\i - 1)+\margin}:{360/\n * (\i)-\margin}:\radius);
		}
	\node {\label};	
\end{tikzpicture}
\par\end{center}%
\end{minipage}

\caption{\label{fig:Graphs-representing-(1,2)(3)-1}Graphs representing the
permutation of $\alpha=(1,2)(3,6,4)(5)$ on $\{1,2,3,4,5,6\}$.}
\end{figure}

Let $\alpha$ be a permutation on $\{1,2,3,4,5,6\}$ defined as follows: 

\begin{tabular}{rlc}
$\alpha(1)=5$ &  & $\alpha(4)=6$\tabularnewline
$\alpha(2)=3$ & and & $\alpha(5)=2$\tabularnewline
$\alpha(3)=1$ &  & $\alpha(6)=4$\tabularnewline
\end{tabular}

Write in $\alpha$ in cycle form.
\end{example}
What we label the elements in $A$ is really arbitrary as a permutation
only maps elements to elements, and any meaning the labels might otherwise
have is irrelevant to this mapping. For example, both the sets $\{bob,sue,bill\}$
or $\{a,b,c\}$ can be represent by the set $\{1,2,3\}$ as they each
have three elements. As such we just as well use the set $\{1,2,\dots,n-1,n\}$.
Each permutations on a set $A=\{1,2,\dots,n-1,n\}$ is an element
in the collection of all permutations on $A$. A collection we will
call the \emph{symmetric group of degree n}, denoted $S_{n}$. 

Of course we first need to show that $S_{n}$ is a group. We have
already seen an example of the group operation we will use, function
composition, in example \ref{ex: simple permutations }. Note that
function operation will always be well defined as the permutations
are all bijections. This operation is can easily be performed using
cycle notation. 
\begin{example}
Let $\sigma=(1,2,3)$ and $\gamma=(1,2)(3)$. $\sigma\circ\gamma$
or just $\sigma\gamma$ as we will write is then 
\[
\begin{array}{rrr}
1\rightarrow2\rightarrow3 & \Rightarrow & (1,3)\\
2\rightarrow1\rightarrow2 & \Rightarrow & (2)\\
3\rightarrow3\rightarrow1 & \Rightarrow & (3,1)
\end{array}
\]
The first and last cycles are equivalent so he have $\sigma\gamma=(1,3)(2)$.

\begin{itemize}
\item Note that by writing a permutation in cycle form it can be inferred
that both these function are bijections of the set $\left\{ 1,2,3\right\} $. 
\end{itemize}
\end{example}

\begin{example}
Let $\alpha=(1,2)(3)(4,5)$ and $\beta=(1,5,3)(2,4)$. 
\begin{eqnarray*}
\alpha\beta & = & (1,4)(2,5,3)
\end{eqnarray*}
\end{example}
\begin{xca}
Let $\alpha=(1,2)(3)(4,5)$ and $\beta=(1,5,3)(2,4)$. Find $\beta\alpha$.
\end{xca}
$S_{n}$ is a nonempty set with a well defined operation. Furthermore
it is easy to see that this operation is closed. 

What would the identity element be? 
\[
e=(1)(2)(3)\dots(n)
\]
$(1)$ means that $1\rightarrow1$. So then $\sigma e=e\sigma$ for
any permutation $\sigma$ in $S_{n}$. 

What about inverses? Permutations are all one-to-one, so each will
have an inverse function. Thus $S_{n}$ is a group.
\begin{example}
Find the inverse of $\beta=(1,5,3)(2,4)$. We need $1\rightarrow5\rightarrow1$.
Also we need $3\rightarrow1\rightarrow3$, and $5\rightarrow3\rightarrow5$
so $(5,1,3)\in\beta^{-1}$. Similarly $2\rightarrow4\rightarrow2$
and $4\rightarrow2\rightarrow4$ implies that $(2,4)\in\beta^{-1}$.
Giving, 
\[
\beta^{-1}=(5,1,3)(2,4)=(1,3,5)(2,4)\mbox{.}
\]
It is easy to check that $\beta\beta^{-1}=(1)(2)(3)(4)(5)=\beta^{-1}\beta$.
\end{example}
\begin{xca}
Find the inverse of $\alpha=(1,2)(3)(4,5)$ and $\alpha\beta=(1,4)(2,5,3)$.
\end{xca}
It is easy to find the number of permutations in $S_{n}$. Let $\sigma\in S_{n}$
and $\sigma(x)=y$ represent the permutation mapping $x\rightarrow y$.
There are $n$ choices for $\sigma(1)$, then $n-1$ choices for $\sigma(2)$,
next there are $n-2$ choices for $\sigma(3)$ etc. So there will
be a total of $n(n-1)(n-2)\cdots1=n!$ permutations in $S_{n}$, i.e.,
$\left|S_{n}\right|=n!$.
\begin{xca}
Find the order of $S_{4}$, $S_{1}$, and $S_{3}$.
\end{xca}

\begin{xca}
Show by example that $S_{3}$ is not abelian. Thus in general symmetric
groups are not abelian. 
\end{xca}

\section{Cosets}

Recall that the elements in $\mathbb{Z}$ each fell into $n$ classes
of $\mathbb{Z}_{n}$ denoted $\left[a\right]$ for an $a\in\mathbb{Z}_{6}$.
For $\mathbb{Z}_{6}$, $2\equiv8\bmod6$ so then $8\in\left[2\right]$.
We can easily find all the elements in $\left[2\right]$ by repeatedly
adding 6 (since we are working with modulo $6$). Similarly we could
find all the elements of each class. 
\begin{eqnarray*}
\left[2\right] & = & \left\{ \dots,-2,0,2,8,14,\dots\right\} \\
\left[3\right] & = & \left\{ \dots,-3,0,3,6,15,\dots\right\} \\
\vdots\\
\left[a\right] & = & \left\{ \dots,a-n,a,a+n,a+2n,\dots\right\} \\
\vdots\\
\left[6\right] & = & \left\{ \dots,-6,0,6,12,\dots\right\} 
\end{eqnarray*}
The elements of each class of $\left[a\right]$ can each be expressed
as $a+kn$ for some $k\in\mathbb{Z}$. It thus would make sense to
write $a+6\mathbb{Z}$ instead of $\left[a\right]$. 

We call sets represented like $3+6\mathbb{Z}=\left\{ \dots,-3,0,3,6,15,\dots\right\} $
\emph{cosets} (note that $6\mathbb{Z}$ is a subgroup of $\mathbb{Z}$).
The idea was first explored by Galois and is key to proving Lagrange's
theorem (the subject of the next section).
\begin{defn}
Let $H$ be a subgroup of $G$ and $a\in G$. 

The set $a\oplus H=\left\{ a\oplus h\colon h\in H\right\} $ is called
the \emph{left coset of H in G containing a}.

The set $H\oplus a=\left\{ h\oplus a\colon h\in H\right\} $ is called
the \emph{right coset of H in G containing a}.
\end{defn}
\begin{example}
\label{ex: Z8  and <2> }Let $G=\mathbb{Z}_{8}$ and $H=\left\langle 2\right\rangle =\left\{ 0,2,4,6\right\} $.
Then 
\begin{eqnarray*}
3\oplus H=3\oplus\left\langle 2\right\rangle  & = & \left\{ 3+0,3+2,3+4,3+6\right\} \\
 & = & \{3,5,7,1\}
\end{eqnarray*}
also since $\mathbb{Z}_{8}$ is abelian, 
\begin{eqnarray*}
H\oplus3=\left\langle 2\right\rangle \oplus3 & = & \left\{ 0+3,2+3,4+3,6+3\right\} \\
 & = & \{3,5,7,1\}
\end{eqnarray*}

The order of $2\in\left\langle 2\right\rangle $, written $\left|2\right|$,
is $4$ since $2+2+2+2=2\cdot4=8\equiv0\bmod8$. $\left|\mathbb{Z}_{8}\right|=8$,
and $2\mid8$, that is $\left|2\right|\mid\left|\mathbb{Z}_{8}\right|$. 
\begin{itemize}
\item We will see that the order of a subgroup (or element) always divides
the order of its group.
\end{itemize}
$4,6,0\in\left\langle 2\right\rangle $ and $4\oplus\left\langle 2\right\rangle =\left\{ 4,6,0,2\right\} =2\oplus\left\langle 2\right\rangle $;
$6\oplus\left\langle 2\right\rangle =\left\{ 4,6,0,2\right\} =2\oplus\left\langle 2\right\rangle $;
and $0\oplus\left\langle 2\right\rangle =\left\langle 2\right\rangle =\left\{ 4,6,0,2\right\} =2\oplus\left\langle 2\right\rangle $\@.
From above $3\notin\{0,2,4,6\}$ and $3\oplus\left\langle 2\right\rangle \neq2\oplus\left\langle 2\right\rangle $. 
\begin{itemize}
\item We will see that cosets are either mutually distinct or equal. As
all elements in $\mathbb{Z}_{8}$ are in $\{0,2,4,6\}$ or $\{1,3,5,7\}$
these are the only \emph{distinct }cosets. Meaning that any other
coset $a\oplus\left\langle 2\right\rangle $ will be congruent to
one of these two cosets. 
\end{itemize}
\end{example}
\begin{xca}
Find all the distinct cosets of $\mathbb{Z}_{9}$.
\end{xca}

\begin{xca}
Consider the group $S_{3}=\{(1),(1,2),(1,3),(2,3)\}$ and its subgroup
$H=\{(1),(1,3)\}$. The distinct left cosets are 

\begin{eqnarray*}
(1)H & = & \left\{ (1)(1),(1)(1,3)\right\} =H\\
(1,2)H & = & \left\{ (1,2)(1),(1,2)(1,3)\right\} =\left\{ (1,2),(1,3,2)\right\} =(1,3,2)H\\
(1,3)H & = & \left\{ (1,3)(1),(1,3)(1,3)\right\} =\left\{ (1,3),(1,3)\right\} =H\\
(2,3)H & = & \left\{ (2,3)(1),(2,3)(1,3)\right\} =\left\{ (2,3),(1,2,3)\right\} =\left(1,2,3\right)H
\end{eqnarray*}
\end{xca}
From these examples/exercises we can see that cosets have some distinct
properties.
\begin{lem}
\label{lem:a is in a+H}Let $H$ be a subgroup of $G$ and $a\in H$.
Then $a\in a\oplus H$ and $a\in H\oplus a$.

\begin{proof}
Let $a\in H$. Since $H$ is a subgroup it has an identity element
$0_{G}$, and $a=a\oplus0_{G}\in a\oplus H$ and $a=0_{G}\oplus a\in a\oplus H$.
\end{proof}
\end{lem}

\begin{lem}
\label{lem:a+H=00003DH iff a in H}Let $H$ be a subgroup of $G$.
If $a\in G$ then $a\oplus H=H$ and $H\oplus a=H$ if and only if
$a\in H$.

\begin{proof}
First we show the result for $a\oplus H=H$ .

$(\Rightarrow)$ Suppose that $a\oplus H=H$. $H$ is a group with
the same identity element as $G$, say $0$. $a\oplus0=a\Rightarrow a\in a\oplus H$
as needed. 

$(\Leftarrow)$ Suppose that $a\in H$ by the lemma above, $a\in a\oplus H$.
As $H$ is close, $a\oplus h\in H$ for any $h\in H$. So then $a\oplus H\subseteq H$. 

Now let $h\in H$. Since $a$ is also in $H$ it follows that $a^{-1}\in H$
and thus $a^{-1}\oplus h\in H$. $h=0\oplus h=(a\oplus a^{-1})\oplus h=a\oplus(a^{-1}\oplus h)\in a\oplus H$.
Thus $H\subseteq a\oplus H$ and it follows that $H=a\oplus H$.

The case for $H\oplus a=H$ follows similarly. 
\end{proof}
\end{lem}
\begin{itemize}
\item The results below are all true for right hand cosets and left hand
cosets. For convenience they are stated for left hand cosets.
\end{itemize}
\begin{thm}
\label{thm:cosets =00003D or disjoint}Let $H$ be a subgroup of $G$.
If $a,b\in G$ then:

\begin{itemize}
\item []$a\oplus H=b\oplus H$ or
\item []$(a\oplus H)\cap(b\oplus H)=\emptyset$. 
\end{itemize}
\begin{proof}
Suppose that the cosets are not disjoint, i.e., $(a\oplus H)\cap(b\oplus H)\neq\emptyset$.
So there is some $x\in(a\oplus H)\cap(b\oplus H)$. Thus there exists
an $h_{1},h_{2}\in H$ such that $x=a\oplus h_{1}$ and $x=b\oplus h_{2}$.
Thus 
\begin{eqnarray*}
a\oplus h_{1} & = & b\oplus h_{2}\\
a & = & (b\oplus h_{2})\oplus h_{1}^{-1}\\
a & = & b\oplus(h_{2}\oplus h_{1}^{-1})
\end{eqnarray*}
Which implies that $a\oplus H=\left(b\oplus(h_{2}\oplus h_{1}^{-1})\right)\oplus H=b\oplus(h_{2}\oplus h_{1}^{-1})\oplus H=b\oplus H$.
As needed. The last equality follows since by closure $h_{2}\oplus h_{1}^{-1}\in H\iff(h_{2}\oplus h_{1}^{-1})\oplus H=H$
by lemma \ref{lem:a+H=00003DH iff a in H}.
\end{proof}
\end{thm}
\begin{cor}
Let $H$ be a subgroup of $G$ and $a,b\in G$. $a\oplus H=b\oplus H$
if and only if $a^{-1}\oplus b\in H$. 

\begin{proof}
$a\oplus H=b\oplus H\iff H=\left(a^{-1}\oplus b\right)\oplus H$ which
implies that $a^{-1}\oplus b\in H$ by the theorem above. 
\end{proof}
\end{cor}
\begin{thm}
\label{thm:|b+H| =00003D |a+H|}Let $H$ be a subgroup of $G$ and
$a,b\in G$. $\left|a\oplus H\right|=\left|b\oplus H\right|$

\begin{proof}
One can show that two sets have the same cardinality by finding a
one-to-one map between them. Consider the map $f\colon a\oplus H\rightarrow b\oplus H$
defined by $f(a\oplus h)=b\oplus h$. By the cancellation property,
theorem \ref{thm:Cancellation-Law.}, $f(a\oplus h_{1})=f(a\oplus h_{2})\Rightarrow b\oplus h_{1}=b\oplus h_{2}\Rightarrow h_{1}=h_{2}$
so $f$ is one-to-one as needed. 
\end{proof}
\end{thm}
\begin{cor}
Let $H$ be a subgroup of $G$ and $a\in G$. $\left|a\oplus H\right|=\left|H\right|$

\begin{proof}
This follows immediately from theorem \ref{thm:|b+H| =00003D |a+H|}
above, since $0_{G}\in G$. So then $\left|a\oplus H\right|=\left|0+H\right|=\left|H\right|$.
\end{proof}
\end{cor}

\begin{thm}
Let $H$ be a subgroup of $G$ and $a\in G$. $a\oplus H$ is a subgroup
of $G$ if and only if $a\in H$.

\begin{proof}
$(\Rightarrow)$ Let $a\oplus H\leq G$. So then $0_{G}\in a\oplus H$.
$0_{G}$ is also in $0_{G}\oplus H=H$ thus $(a\oplus H)\cap H\neq\emptyset$
and it follows that $a\oplus H=H$ by theorem \ref{thm:cosets =00003D or disjoint}. 

$(\Leftarrow)$ Suppose that $a\in H$. This implies that $a\oplus H=H$
by lemma \ref{lem:a+H=00003DH iff a in H}, and $a\oplus H=H\leq G$. 
\end{proof}
\end{thm}

\section{Lagrange's Theorem}

Our study of group theory ends with Lagrange's theorem and its applications.
Lagrange's theorem is the most important result in finite group theory.
It provides a necessary characterization for the order of subgroups
(and consequently elements) of a finite group. Recall example \ref{ex: Z8  and <2> }: 
\begin{example}
$\mathbb{Z}_{8}$ and $\left\langle 2\right\rangle =\left\{ 0,2,4,6\right\} $.

\begin{itemize}
\item $\left\langle 2\right\rangle \leq\mathbb{Z}_{8}$: Since $2\in\mathbb{Z}_{8}$,
$\left\langle 2\right\rangle $ is the cyclic subgroup generated by
$2$. Recall that it is also abelian as all cyclic groups are abelian
(theorem \ref{thm:All-cyclic-groups-are-abelian}).
\item $\left|2\right|=\left|\left\langle 2\right\rangle \right|=4$: Since
$2^{4}=2+2+2+2=8\equiv0\bmod8$. 
\item Note that $4\mid8$: This is no coincidence. We will see that Lagrange's
theorem tell us that this is a necessary condition of all subgroups
and group elements.
\item Recall that are cosets of $\left\langle 2\right\rangle $ are either
distinct or disjoint, e.g., $g_{1}+\left\langle 2\right\rangle =g_{2}+\left\langle 2\right\rangle $
or $\left(g_{1}+\left\langle 2\right\rangle \right)\cap\left(g_{2}+\left\langle 2\right\rangle \right)=\emptyset$.
So the distinct left/right cosets of a subgroup \emph{partition }its
group. This will be important in the proof of Lagrange's theorem.
\end{itemize}
\end{example}
\begin{lem}
\label{lem:G is the union of cosets}If $H$ is a subgroup of a finite
group $G$ then $G$ is the union of the distinct cosets of $H$,
i.e., 
\[
G=a_{1}\oplus H\cup a_{2}\oplus H\cup\cdots\cup a_{k}\oplus H
\]
 for a list of elements $a_{1},a_{2},\dots,a_{k}\in G$ where $a_{i}\oplus H\neq a_{j}\oplus H$
if $i\neq j$ for any $1\leq i,j\leq k\leq$. 

\begin{proof}
Let $a_{1}$ be any element in $G$. If there is no such element $a_{2}\in G$
such that $a_{1}\oplus H\neq a_{2}\oplus H$ then $a_{2}\in a_{1}\oplus H$
for any $a_{2}\in G$. Which implies that $a_{1}\oplus H=G$, and
we are done. 

Otherwise we can either choose another $a_{2}$ such that $a_{1}\oplus H\neq a_{2}\oplus H$.
As above, either we can find another element $a_{3}\in G$ such that
$a_{3}\oplus H$ will be distinct from $a_{1}\oplus H$ and $a_{2}\oplus H$
or $a_{1}\oplus H\cup a_{2}\oplus H=G$. $G$ is finite, so continuing
like this we can compile a list of distinct elements in $G$, $a_{1},a_{2},\dots,a_{k}$
for $k\leq\left|G\right|$, such that $G=a_{1}\oplus H\cup a_{2}\oplus H\cup\cdots\cup a_{k}\oplus H$.
\end{proof}
\end{lem}
While the cosets themselves will be unique the list of elements $a_{1},a_{2},\dots,a_{k}$
is not. However the number of cosets $k$ will always be the same.
This number is important enough to have a name -the \emph{index }of
the subgroup $H$.
\begin{defn}
If $H$ is a subgroup of a finite group $G$ such that $G=a_{1}\oplus H\cup a_{2}\oplus H\cup\cdots\cup a_{k}\oplus H$
then the $k$ is called the \emph{index }of $H$ in $G$; denoted
$\left|G\colon H\right|$. As we will see below, $\left|G\colon H\right|=\nicefrac{\left|G\right|}{\left|H\right|}$
so the index is sometimes notated $\left|G\right|/\left|H\right|$.
\end{defn}
\begin{thm}
\textbf{\emph{Lagrange's Theorem. }}If $H$ is a subgroup of a finite
group $G$ then $\left|H\right|\mid\left|G\right|$. 

\begin{proof}
Let $a_{1}\oplus H,a_{2}\oplus H,\dots,a_{k}\oplus H$ be the distinct
left cosets of $H$ in $G$. By lemma \ref{lem:G is the union of cosets}
above, 
\[
G=a_{1}\oplus H\cup a_{2}\oplus H\cup\cdots\cup a_{k}\oplus H
\]
. Since this union is disjoint we have, 
\[
\left|G\right|=\left|a_{1}\oplus H\right|+\left|a_{2}\oplus H\right|+\cdots+\left|a_{k}\oplus H\right|\mbox{.}
\]
By the corollary to theorem \ref{thm:|b+H| =00003D |a+H|}, each coset
$a_{i}\oplus H$ in this union has exactly $\left|H\right|$ elements,
i.e., $\left|a_{1}\oplus H\right|=\left|H\right|$. So then $\left|G\right|=k\cdot\left|H\right|$.
\end{proof}
\end{thm}
\begin{example}
Consider $\mathbb{Z}_{8}$ and its subgroup $\left\langle 2\right\rangle $.
The order of $\left\langle 2\right\rangle $ is either $1,2,4$, or
$8$ as these are the only divisors of $\left|\mathbb{Z}_{8}\right|=8$
. 
\end{example}

\begin{example}
$H=\left\{ 0,1,2\right\} $ is not a subgroup of $\mathbb{Z}_{8}$
as $\left|H\right|=3$ and $3\nmid8$.
\end{example}

\begin{example}
Let $G$ be a group and $H\leq G$. If $\left|G\right|=35=7\cdot5$.
The order of $H$ is $1,5,7,$ or $35$.
\end{example}

\begin{example}
Let $G$ be a group and $H\leq G$. If $\left|G\right|=40$ and $\left|H\right|=7$.
Then $H\nleq G$.
\end{example}
Note: The converse of Lagrange's theorem is false. Simply because
the order of a subset divides the order of a group \uline{does not
make it a subgroup}.
\begin{example}
(The converse of Lagrange's theorem is false). Let $H=\{1,2,3,4\}\subseteq\mathbb{Z}_{8}$.
$\left|H\right|=4$ and $4\mid8$, but $H\nleq G$ as $0\notin G$.
So it has no identity element.
\end{example}

\begin{cor}
If $G$ is a group and $a\in G$ then $\left|a\right|\mid\left|G\right|$.

\begin{proof}
Recall that $\left|\left\langle a\right\rangle \right|=\left|a\right|$
for any element $a$ in a group $G$, and that $\left\langle a\right\rangle \leq G$.
So as an immediate consequence of Lagrange's theorem $\left|a\right|\mid\left|G\right|$.
\end{proof}
\end{cor}

\begin{cor}
If $G$ is a group with prime order then $G$ is cyclic.

\begin{proof}
As $\left|G\right|>1$, we can choose an element $a\in G$ such that
$a\neq0_{G}$. Since $\left|G\right|=p$ for some prime integer $p$,
$\left|a\right|\mid p$ by Lagrange's theorem. Thus $\left|a\right|=1$
or $p$, but $\left|a\right|\neq1$ since $a\neq0_{G}$. It follows
that $\left|\left\langle a\right\rangle \right|=p$, and since $\left\langle a\right\rangle \subseteq G$
the result follows. 
\end{proof}
\end{cor}

\end{document}
